\section{Quantitative liveness: expected termination under every scheduler}
\label{sec:bellman}

\subsection{A graph cycle need not prevent probabilistic termination}

A retry loop that succeeds with probability $1/4$ on each iteration has a
self-loop. There is an infinite path consisting entirely of retries, so no
natural-valued function can strictly decrease on every possible transition.
Nevertheless, the probability of that infinite path is zero and the expected
number of iterations is 4. Forge's existing pathwise ranking mechanisms cannot
certify this merely by searching for more elaborate deterministic ranks.
The needed local operator is expectation.

We consider a finite Markov decision process (MDP). Every state $s$ has at least
one action; action $a$ induces a rational distribution $P_{s,a}$ over states.
The designated terminal states are absorbing in the input itself. The stopping
time $T$ is the first time a terminal state is entered, with $T=\infty$ if
this never happens. Each nonterminal step has cost 1. A scheduler may depend
on the complete finite history and may randomize among available actions.
The goal is an upper bound valid for \emph{every} such scheduler.

The terminal and cost conventions matter. A zero-cost loop can have expected
cost zero while never terminating. A checker that omits the unit cost from its
inequality would prove the wrong property. Likewise, rewriting a nonabsorbing
state into a terminal self-loop changes the model unless the source stopping
semantics explicitly justifies that transformation.

\subsection{A local Bellman certificate}

A positive certificate supplies a state region $W$ containing the initial
state, closed under every action, and a rational function $V:W\to\Q_{\ge0}$.
The values at terminal states are zero. For each nonterminal $s\in W$ and
\emph{every} available action $a$, the checker requires
\begin{equation}
 V(s)\ \ge\ 1+\sum_tP_{s,a}(t)V(t).
 \label{eq:bellman-upper}
\end{equation}
Closure ensures that no mass escapes to an unrepresented potential value.
The delivered certificate also supplies a deterministic positional policy
$\pi$ attaining equality at every nonterminal state in $W$.

\begin{theorem}[Uniform expected-time upper bound]
\label{thm:bellman-upper}
If \eqref{eq:bellman-upper} holds on a universally closed region containing
$s_0$, then for every history-dependent randomized scheduler,
\[
 \E_{s_0}[T]\le V(s_0),\qquad \Prob_{s_0}(T=\infty)=0.
\]
For every positive integer $N$, it also gives
$\Prob_{s_0}(T>N)\le V(s_0)/N$.
\end{theorem}
\begin{proof}
Condition on a finite execution history that has not yet terminated.
The scheduler chooses a distribution over actions. Since
\eqref{eq:bellman-upper} holds for each action, it also holds for their
convex combination. Thus the conditional expected change in the potential is
at most $-1$. Once terminal, both the additional cost and the potential change
are zero.

Apply this inequality for exactly $N$ steps to the stopped process. Finite
induction and iterated expectation give
\begin{equation}
 \E[\min(T,N)] + \E[V(X_{\min(T,N)})]\le V(s_0).
 \label{eq:stopped-bound}
\end{equation}
Dropping the nonnegative second term and increasing $N$ proves the expected-time
bound by monotone convergence. A positive mass at $T=\infty$ would force
$\E[\min(T,N)]$ to grow without bound, so that mass is zero. Finally
$N\Prob(T>N)\le\E[T]$ gives the tail estimate.
\end{proof}

The proof uses a finite-horizon inequality before taking a limit. A Lean
implementation need not begin by invoking an optional-stopping theorem whose
integrability side conditions it has not yet established. It can first define
finite execution-tree expectations by recursion and prove
\eqref{eq:stopped-bound}. The measure-theoretic interpretation is a separate,
reusable bridge.

\begin{theorem}[Exact worst-case expectation]
\label{thm:bellman-tight}
If a positional policy $\pi$ makes \eqref{eq:bellman-upper} an equality at
every nonterminal state of $W$, then
\[
 \sup_{\text{all schedulers }\sigma}\E^\sigma_{s_0}[T]
   =\E^\pi_{s_0}[T]=V(s_0).
\]
\end{theorem}
\begin{proof}
For policy $\pi$, the finite telescoping argument yields equality in
\eqref{eq:stopped-bound}. The upper-bound theorem already establishes almost-sure
termination under this policy. Since $W$ is finite, $V$ is bounded by a finite
constant $M$, and the remaining-potential term is at most
$M\Prob(T>N)$, which tends to zero. The equality therefore converges to
$\E^\pi[T]=V(s_0)$. The upper-bound theorem bounds all other schedulers by
the same number.
\end{proof}

A certificate intended only to prove a budget could omit the tight policy and
accept any valid super-solution. The prototype uses the stronger exact-optimum
contract to make differential testing against exhaustive policy evaluation
possible. A budget query compares $V(s_0)$ to the original budget only after
this certificate is checked.

\subsection{A finite certificate of nontermination}

A negative certificate supplies a nonempty set $U$ of nonterminal states and
one action $\rho(s)$ at each $s\in U$, such that
\begin{equation}
 \supp(P_{s,\rho(s)})\subseteq U.
 \label{eq:trap}
\end{equation}
It also supplies a finite path from $s_0$ to $U$, labeling each edge with a
legal action and a strictly positive transition probability. The checker
multiplies these probabilities and verifies the claimed positive product
$\alpha$.

\begin{proposition}[Reachable-trap soundness]
\label{prop:trap}
A valid path-and-trap certificate establishes a scheduler with
$\Prob(T=\infty)\ge\alpha>0$, hence infinite expected termination time.
\end{proposition}
\begin{proof}
Use a history-dependent scheduler that tries the listed path and, when that
path has been followed, switches to $\rho$ in $U$. The probability of following
the entire finite path is $\alpha$. Conditional on that event, all subsequent
transitions stay in $U$ with probability one and no terminal state is reached.
A nonnegative random variable that is infinite on a positive-probability event
has infinite expectation.
\end{proof}

The witness scheduler may use finite memory during the prefix; it need not be
a single stationary table compatible with every occurrence along the path.
The universal upper-bound goal allows all history-dependent schedulers, so
this is a valid counterexample. An unreachable nonterminal trap is not a
counterexample from $s_0$, which is why the path cannot be omitted.

\subsection{A complete algorithm on the stated finite fragment}

First compute $W$, the states reachable from $s_0$ under the union of all
positive-probability action edges. Let $U_0=W\setminus\mathrm{Terminal}$ and
iterate
\begin{equation}
 U_{k+1}=\{s\in U_k:\exists a,\ \supp(P_{s,a})\subseteq U_k\}.
 \label{eq:trap-fp}
\end{equation}
The sequence decreases and stabilizes after at most $|W|$ strict rounds.
If its limit is nonempty, select one staying action at each surviving state
and recover a reachability path from the initial search. This yields
Proposition~\ref{prop:trap}.

Suppose instead that the limit is empty. Assign each nonterminal state the
round in which it is removed. At removal, \emph{every} action has some positive
probability of moving to a terminal state or a state removed earlier. Let
$\eta>0$ be the smallest positive transition probability in the reachable
finite model. Under any scheduler and after any history, there is probability
at least $\eta^{|W|}$ of following rank-decreasing successors to termination
within $|W|$ steps. Repeating the argument in blocks gives a uniform geometric
tail and a finite expected-time bound. This establishes that all stationary
policies are proper before any policy matrix is inverted.

Now perform maximizing policy iteration. For a deterministic stationary policy
$\pi$, let $Q_\pi$ be its transition matrix restricted to reachable nonterminal
states. Evaluate it exactly by
\begin{equation}
 (I-Q_\pi)v_\pi=\boldsymbol 1.
 \label{eq:policy-eval}
\end{equation}
At each state, compute $1+P_{s,a}\cdot v_\pi$ for all actions. Switch only when
an action is strictly greater than the current value. Re-evaluate and repeat.
The prototype uses exact Gauss--Jordan elimination, not rounded matrix
inversion.

\begin{lemma}[Policy improvement after the trap test]
\label{lem:policy-improve}
When the trap fixed point is empty, every policy evaluation is nonsingular,
and a strictly improving switch cannot cause a policy cycle.
\end{lemma}
\begin{proof}
Properness and the uniform geometric tail imply
$(I-Q_\pi)^{-1}=\sum_{k\ge0}Q_\pi^k$, a finite-entry nonnegative matrix.
For a switched policy $\pi'$, let
$d=\boldsymbol1+Q_{\pi'}v_\pi-v_\pi$. The switching rule makes $d\ge0$ and
at least one component positive. Then
\[
 v_{\pi'}-v_\pi=(I-Q_{\pi'})^{-1}d\ge0.
\]
Because the inverse includes its identity summand, at least one component
strictly increases. A return to an earlier policy would return to its unique
value vector, contradicting monotone strict progress. There are finitely many
stationary policies, so iteration terminates.
\end{proof}

At termination, the current policy has equality in its evaluation equations and
no action exceeds its value. The vector and policy therefore satisfy both
Bellman certificate checks.

\begin{theorem}[Finite-fragment decision result]
\label{thm:mdp-complete}
For a finite total MDP with exact rational probabilities and absorbing terminal
states, the algorithm terminates with either an exact worst-case expected-time
certificate or a reachable nonterminal trap certificate. The former proves
finite expected time under every scheduler; the latter refutes that property.
\end{theorem}
\begin{proof}
The finite trap iteration terminates. A nonempty result gives the negative
certificate. An empty result establishes uniform properness, after which
Lemma~\ref{lem:policy-improve} guarantees termination of policy iteration.
The final local equations give the positive certificate and its correctness by
Theorems~\ref{thm:bellman-upper} and~\ref{thm:bellman-tight}.
\end{proof}

The dichotomy relies essentially on finiteness. It does not assert that every
almost-surely terminating infinite-state program has finite expected time. Nor
does failure to find a symbolic polynomial potential refute termination.

\subsection{Worked scheduler examples}

\paragraph{Geometric retry with an adversary.}
There are two states: retry $r$ and terminal $z$. At $r$, an adversary chooses
between success probabilities $1/2$ and $1/4$; unsuccessful trials return to
$r$. Let $V(r)=4$ and $V(z)=0$. The two Bellman right-hand sides are
\[
 1+\tfrac12\cdot4=3,\qquad
 1+\tfrac34\cdot4=4.
\]
Thus every scheduler has expected duration at most 4, and always selecting the
slower action attains 4. No graph edge has been removed: the retry self-loop
is present in both the model and the proof.

\paragraph{A good policy does not prove a universal bound.}
At $r$, one action terminates immediately and another stays at $r$ with
probability one. An existential controller can terminate in one step, but the
universal scheduler statement is false. The trap $\{r\}$ with the staying
action and the empty prefix has $\alpha=1$. The prototype checks that it returns
this negative outcome rather than the expected duration of the convenient
terminating policy.

\paragraph{An unreachable trap is irrelevant.}
Add a separate absorbing nonterminal state $u$ with no path from $s_0$. The
reachability slice removes it before the trap test. A hand-written negative
certificate naming $u$ without a valid path is rejected. This demonstrates a
necessary source of root dependence: changing the initial state can change
the theorem even when every transition table remains identical.

\subsection{A symbolic extension that reuses existing Forge machinery}

For an infinite-state probabilistic loop, let the update be $x'=F(x,\xi)$
with a finitely supported random choice $\xi$. Given a polynomial potential
$V$, the local drift obligation is
\[
 V(x)-\sum_{\xi}p_\xi V(F(x,\xi))\ge\varepsilon>0
\]
on the nonterminal invariant, together with $V\ge0$. Forge's existing exact
polynomial and cone workers can discharge this obligation once the expected
substitution is reified. They should not demand the stronger and often false
condition $V(x)-V(F(x,\xi))\ge\varepsilon$ for each outcome.

For example, a walk on positive integers moves down with probability $p$ and
up with probability $q=1-p$, with $p>q$, and stops at zero. Taking
$V(n)=n/(p-q)$ gives expected drift exactly 1 at every positive state. The
finite-horizon telescoping argument gives $\E[T]\le V(n)$ without a pointwise
rank. This is a designed composition with existing Forge certificates, not a
symbolic-search feature executed by the delivered finite-state prototype.
State-dependent probabilities require their own nonnegativity and normalization
proofs; clearing a denominator without proving its sign is not an acceptable
translation.
